\hypertarget{introduction}{%
\section{Introduction}\label{introduction}}

One-shot federated distillation constructs a global student from independently trained local models without further rounds of client optimization. Distillation transfers an ensemble's outputs into a single model \cite{hinton2015}, and federated model fusion extends this principle to client models \cite{li2019,lin2020}. Under heterogeneous class allocations, however, the server must decide which teachers should contribute to each proxy example and which parts of their predictions should be transferred.

Class presence and predictive competence provide different answers to the first question. A teacher may have encountered a class without predicting it reliably. Conversely, excluding teachers that never trained on a class may already explain much of the improvement obtained by a more elaborate competence rule. A comparison against uniform aggregation alone cannot distinguish these explanations. This distinction is particularly relevant when classes are distributed in disjoint groups: a competence mask may reproduce the training-presence mask exactly.

Contributor selection and output restriction address different questions. After choosing a teacher for a proxy example, the server can preserve its complete class distribution or remove probabilities outside its accredited classes. Complete distributions may convey useful interclass relationships, but restriction also changes their concentration. Moreover, if selection uses the true proxy label, restriction can improve target likelihood by construction. An apparently better target is therefore not sufficient evidence of a better distilled student.

We study these choices with a labeled public proxy and a binary expertise mask estimated after teacher checkpoint selection. EXPERT selects teachers accredited for the proxy label and averages their full probability vectors. A training-presence control replaces accreditation with a record of observed training classes, while ORACLE selects teachers whose top prediction is correct for the individual sample. We compare pooling operators and restrict the outputs of the same selected experts. Related studies already consider certainty-weighted federated distillation and monoclass teachers \cite{sattler2021,maron2025}; our question concerns the incremental utility and limitations of these specific information sources under a shared experimental protocol.

The availability of public labels also changes the practical reference point. A student can learn directly from those labels without accessing teachers. We therefore compare EXPERT and uniform distillation with proxy-only cross-entropy training over nested proxy subsets. This separates recovering performance relative to an unsuitable ensemble target from adding value beyond direct supervision.

Our contributions are threefold. First, we evaluate training presence and measured competence as distinct routing signals using the same teachers and student-training controls. Second, we separate contributor selection, pooling and output restriction, and show why their target-level diagnostics must be interpreted alongside student accuracy and likelihood. Third, we evaluate both uniform and expertise-based distillation against direct supervision across public-label budgets. The contribution is a controlled empirical characterization of these choices, not a claim that teacher selection or specialist distillation is itself new.
